% Chapter: PEG校准与季节性利润校准
\section{PEG/PSG校准排名}

\begin{table}[H]
\centering\scriptsize
\renewcommand{\arraystretch}{1.1}
\begin{tabularx}{\textwidth}{c l c c c c l}
\toprule
\textbf{排名} & \textbf{标的} & \textbf{2026E PE} & \textbf{增速} & \textbf{PEG} & \textbf{结论} & \textbf{两套体系对照} \\
\midrule
1 & 协创数据 & 38$\times$ & +180--230\% & \textbf{0.17--0.21} & \textcolor{riskgreen}{极度低估} & PEG低估+机构B级=首选 \\
2 & 工业富联 & 19$\times$ & +50\% & \textbf{0.32--0.48} & \textcolor{riskgreen}{低估} & PEG低估+机构A+=防御首选 \\
3 & 中际旭创 & 29$\times$ & +40--70\% & 0.41--0.73 & \textcolor{riskgreen}{低估} & 算力底座双雄 \\
4 & 新易盛 & 32$\times$ & +50--70\% & 0.45--0.64 & \textcolor{riskgreen}{低估} & 光模块次龙头 \\
5 & 智微智能 & 57$\times$(季校46$\times$) & +150--200\% & 0.29--0.38 & \textcolor{blue}{PEG低但解禁} & 9月首发解禁→机构禁买 \\
6 & 奥比中光 & 160$\times$(季校144$\times$) & +200--300\% & 0.53--0.80 & \textcolor{riskamber}{分歧最大} & PEG似低估但机构D级 \\
7 & 天准科技 & 扭亏(PS视角) & 营收+60--80\% & PSG 0.13--0.17 & \textcolor{blue}{PSG低估} & 7月解禁后重估 \\
8 & 柯力传感 & 49$\times$ & +25--60\% & 0.82--1.96 & 中性 & 等Optimus投产催化 \\
9 & 德赛西威 & 21$\times$ & +5--15\% & 1.5--4.2 & \textcolor{blue}{远期PEG骤降} & 2027E PEG 0.37=首选 \\
10 & 绿的谐波 & 397$\times$(季校133$\times$) & +50--70\% & \textbf{5.7--7.9} & \textcolor{riskred}{全表最贵} & 无论怎么调PEG$>$2 \\
11 & 鸣志电器 & 222$\times$(季校298$\times$) & -20\textasciitilde+25\% & $>$8.9 & \textcolor{riskred}{极端泡沫} & PEG框架无效 \\
\bottomrule
\end{tabularx}
\caption{PEG视角估值排名（含季节性校准后PE）}
\end{table}

\section{季节性利润分布}

\begin{table}[H]
\centering\small
\renewcommand{\arraystretch}{1.2}
\begin{tabularx}{\textwidth}{l c c c c X}
\toprule
\textbf{子板块} & \textbf{Q1占比} & \textbf{Q2} & \textbf{Q3} & \textbf{Q4} & \textbf{关键驱动} \\
\midrule
人形机器人零部件 & 15\% & 22\% & 30\% & \textbf{33\%} & Q3-Q4备货高峰+年底回款；Q1春节停产 \\
智驾Tier1/域控 & 18\% & 25\% & 28\% & \textbf{29\%} & Q4车企冲量；Q1春节+价格战影响最重 \\
AI算力/服务器 & 27\% & 25\% & 24\% & 24\% & 基本均匀；云厂商Capex波动小 \\
光模块/通信硬件 & 23\% & \textbf{27\%} & 25\% & 25\% & Q2海外下单高峰 \\
\bottomrule
\end{tabularx}
\caption{各子板块季度利润分布系数}
\end{table}

\section{季节性校准对估值的影响}

\begin{table}[H]
\centering\scriptsize
\renewcommand{\arraystretch}{1.1}
\begin{tabularx}{\textwidth}{l c c c c c X}
\toprule
\textbf{标的} & \textbf{Q1净利} & \textbf{Q1$\times$4} & \textbf{Q1占比} & \textbf{校准全年} & \textbf{校准PE} & \textbf{结论} \\
\midrule
协创数据 & 7.5亿 & 30亿 & 27\% & 27.8亿 & 37$\times$ & 与一致预期吻合 \\
德赛西威 & 4.61亿 & 18.4亿 & 18\% & \textbf{25.6亿} & \textbf{21$\times$} & 全年低点确认，H2弹性大 \\
工业富联 & 106亿 & 424亿 & 27\% & 393亿 & 21$\times$ & 算力底座均匀，校准影响小 \\
绿的谐波 & 0.87亿 & 3.48亿 & 15\% & \textbf{5.80亿} & 133$\times$ & 即使校准PE仍133$\times$，结论不变 \\
鸣志电器 & 0.14亿 & 0.56亿 & 15\% & 0.93亿 & 298$\times$ & PE仍近300$\times$，极端泡沫 \\
奥比中光 & 0.31亿 & 1.24亿 & 16\% & 1.94亿 & 144$\times$ & 接近乐观假设上限 \\
智微智能 & 1.09亿 & 4.36亿 & 20\% & \textbf{5.45亿} & 46$\times$ & 校准后PE从57降到46 \\
\bottomrule
\end{tabularx}
\caption{Q1利润→季节性校准全年→校准PE}
\end{table}

\textbf{校准结论}：季节性校准对德赛西威（确认Q1为全年低点）和智微智能（PE显著下降）影响最大；对绿的谐波/鸣志无本质改变（仍为极端泡沫）。

\section{三类H2放大量叙事分类}

\begin{table}[H]
\centering\small
\renewcommand{\arraystretch}{1.2}
\begin{tabularx}{\textwidth}{c X X c}
\toprule
\textbf{类别} & \textbf{特征} & \textbf{代表标的} & \textbf{兑现概率} \\
\midrule
A类（高确定） & Q1已有订单锁定H2产能，季节性放量有历史规律支撑 & 工业富联、中际旭创、新易盛 & 80\%+ \\
B类（中确定） & H2放量依赖特定催化剂落地（产品发布/IPO/定点） & 德赛(Thor量产)、天准(域控)、协创(Coalition) & 50--70\% \\
C类（低确定） & H2放量纯属叙事外推，无订单/产能佐证 & 绿的谐波、鸣志、中大力德、索辰 & $<$30\% \\
\bottomrule
\end{tabularx}
\caption{三类H2放大量叙事分类}
\end{table}

\section{四维校准后最终综合空间}

综合第九节基础估值 $\times$ 传导系数修正 $\times$ 机构可投性修正 $\times$ 季节性校准后的最终结论：

\begin{table}[H]
\centering\scriptsize
\renewcommand{\arraystretch}{1.15}
\begin{tabularx}{\textwidth}{l c c c c c l}
\toprule
\textbf{标的} & \textbf{九节空间} & \textbf{传导修正} & \textbf{机构修正} & \textbf{季节校准} & \textbf{最终空间} & \textbf{操作建议} \\
\midrule
德赛西威 & +62--81\% & 上调+20\% & $\times$1.15 & Q1低点确认 & \textcolor{riskgreen}{\textbf{+107--142\%}} & 首选重仓 \\
协创数据 & -5--+8\% & 上调+15\% & $\times$1.05 & 吻合一致预期 & \textcolor{riskgreen}{\textbf{+21--49\%}} & 低吸加仓 \\
工业富联 & +40--60\% & 确认 & $\times$1.10 & 均匀无影响 & \textcolor{riskgreen}{\textbf{+44--66\%}} & 防御首选 \\
中际旭创 & +25--50\% & 确认 & $\times$1.10 & Q2偏强 & \textcolor{riskgreen}{\textbf{+27--54\%}} & 算力配置 \\
新易盛 & +20--45\% & 确认 & $\times$1.10 & Q2偏强 & \textcolor{riskgreen}{\textbf{+22--50\%}} & 算力配置 \\
天准科技 & 0--+23\% & 上调 & $\times$0.85(解禁) & 扭亏待验 & \textcolor{blue}{+10--35\%} & 7月解禁后 \\
柯力传感 & -11--+14\% & 确认 & $\times$0.95 & Q1低+Optimus & \textcolor{blue}{+5--20\%} & 等7月投产 \\
绿的谐波 & -87---85\% & 不变 & $\times$0.70(D级) & 仍133$\times$ & \textcolor{riskred}{\textbf{-84---80\%}} & 规避 \\
鸣志电器 & -73---82\% & 不变 & $\times$0.70 & 仍298$\times$ & \textcolor{riskred}{\textbf{-84\%}} & 规避 \\
奥比中光 & -60---70\% & 上修(跨阵营) & $\times$0.75 & 接近乐观上限 & \textcolor{riskred}{-44---21\%} & 泡沫区观望 \\
索辰科技 & -50\% & 不变 & $\times$0.65(D-级) & Q4兑现为主 & \textcolor{riskred}{-67\%} & 规避 \\
中大力德 & -45\% & 不变 & $\times$0.70 & H2起量但烂基本面 & \textcolor{riskred}{-60\%} & 规避 \\
\bottomrule
\end{tabularx}
\caption{四维校准后最终综合空间（投资者应优先参考此表）}
\end{table}

\begin{keyinsight}[最终结论]
大资金唯一可重仓5票：\textbf{德赛西威(20\%)、协创数据(15\%)、工业富联(15\%)、中际旭创(12\%)、新易盛(10\%)}，合计建议仓位60--70\%。其余标的单票$\leq$5\%或0持仓。
\end{keyinsight}
